\section{Quasi-affine morphisms; quasi-projective morphisms. Proper
morphisms; projective morphisms}
\label{section:II.5}

\subsection{Quasi-affine morphisms.}
\label{subsection:II.5.1}

\begin{definition}[5.1.1]
\label{II.5.1.1}
We call a scheme isomorphic to a subscheme induced on a quasi-compact open
subset of an affine scheme a \emph{quasi-affine scheme}. We say that a
morphism $f:X\to Y$ is \emph{quasi-affine}, or also that $X$ (considered as a
$Y$-prescheme by means of $f$) is a \emph{quasi-affine $Y$-scheme}, if there
exists a covering $(U_\alpha)$ of $Y$ by affine open subsets such that the
$f^{-1}(U_\alpha)$ are quasi-affine schemes.
\end{definition}

It is clear that a quasi-affine morphism is \emph{separated}
(\hyperref[I.5.5.5]{I, 5.5.5} and \hyperref[I.5.5.8]{5.5.8}) and
\emph{quasi-compact} (\hyperref[I.6.6.1]{I, 6.6.1}); every affine
morphism is evidently quasi-affine.

Recall that for every prescheme $X$, if we set
$A=\Gamma(X,\mathcal{O}_X)$, the identity homomorphism
$A\to A=\Gamma(X,\mathcal{O}_X)$ defines a morphism
$X\to\operatorname{Spec}(A)$ called \emph{canonical}
(\hyperref[I.2.2.4]{I, 2.2.4}); moreover, it is none other than the
canonical morphism defined in (\hyperref[II.4.5.1]{4.5.1}) for the particular
case where $\mathcal{L}=\mathcal{O}_X$, if we recall that
$\operatorname{Proj}(A[T])$ is canonically identified with
$\operatorname{Spec}(A)$ (\hyperref[II.3.1.7]{3.1.7}).

\begin{proposition}[5.1.2]
\label{II.5.1.2}
Let $X$ be a quasi-compact scheme or a prescheme whose underlying space
is Noetherian, and $A$ the ring $\Gamma(X,\mathcal{O}_X)$. The following
conditions are equivalent:
\begin{enumerate}
  \item[\rm{a)}] $X$ is a quasi-affine scheme.
  \item[\rm{b)}] The canonical morphism
  $u:X\to\operatorname{Spec}(A)$ is an open immersion.
  \item[\rm{b')}] The canonical morphism
  $u:X\to\operatorname{Spec}(A)$ is a homeomorphism of $X$ onto a
  subspace of the underlying space of $\operatorname{Spec}(A)$.
  \item[\rm{c)}] The $\mathcal{O}_X$-module $\mathcal{O}_X$ is very ample
  relative to $u$ (\hyperref[II.4.4.2]{4.4.2}).
  \item[\rm{c')}] The $\mathcal{O}_X$-module $\mathcal{O}_X$ is ample
  (\hyperref[II.4.5.1]{4.5.1}).
  \item[\rm{d)}] As $f$ ranges over $A$, the $X_f$ form a base of the
  topology of $X$.
  \item[\rm{d')}] As $f$ ranges over $A$, those $X_f$ which are affine
  form a covering of $X$.

\oldpage[II]{95}

  \item[\rm{e)}] Every quasi-coherent $\mathcal{O}_X$-module is generated by
  its sections over $X$.
  \item[\rm{e')}] Every quasi-coherent sheaf of ideals of
  $\mathcal{O}_X$ of finite type is generated by its sections over $X$.
\end{enumerate}
\end{proposition}

It is clear that b) implies a), and a) implies c) by virtue of criterion b) of
(\hyperref[II.4.4.4]{4.4.4}) applied to the identity morphism (taking into
account the remark preceding the statement); on the other hand, c) implies c')
(\hyperref[II.4.5.10]{4.5.10, (i)}), and c'), b) and b') are equivalent by
(\hyperref[II.4.5.2]{4.5.2, criteria b) and b')}). Finally, c') is equivalent
to each of the criteria d), d'), e), e') by virtue of
(\hyperref[II.4.5.2]{4.5.2, criteria a), a'), c)}) and
(\hyperref[II.4.5.5]{4.5.5, criterion d'')}).

We observe moreover that, with the preceding notation, those $X_f$ which
are affine form a \emph{base} of the topology of $X$, and that the
canonical morphism $u$ is \emph{dominant}
(\hyperref[II.4.5.2]{4.5.2}).

\begin{corollary}[5.1.3]
\label{II.5.1.3}
Let $X$ be a quasi-compact prescheme. If there exists a morphism $v:X\to Y$ of
$X$ into an affine scheme $Y$ which is a homeomorphism of $X$ onto an
open subspace of $Y$, then $X$ is quasi-affine.
\end{corollary}

Indeed, there exists a family $(g_\alpha)$ of sections of $\mathcal{O}_Y$
over $Y$ such that the $D(g_\alpha)$ form a base of the
topology of $v(X)$; if $v=(\psi,\theta)$ and if we set
$f_\alpha=\theta(g_\alpha)$, we have
$X_{f_\alpha}=\psi^{-1}(D(g_\alpha))$
(\hyperref[I.2.2.4.1]{I, 2.2.4.1}), hence the $X_{f_\alpha}$ form a
base of the topology of $X$, and criterion d) of
(\hyperref[II.5.1.2]{5.1.2}) is satisfied.

\begin{corollary}[5.1.4]
\label{II.5.1.4}
If $X$ is a quasi-affine scheme, every invertible $\mathcal{O}_X$-module is
very ample (relative to the canonical morphism) and \emph{a fortiori} ample.
\end{corollary}

Indeed, such a module $\mathcal{L}$ is generated by its sections over
$X$ (\hyperref[II.5.1.2]{5.1.2, e)}), hence
$\mathcal{L}\otimes\mathcal{O}_X=\mathcal{L}$ is very ample
(\hyperref[II.4.4.8]{4.4.8}).

\begin{corollary}[5.1.5]
\label{II.5.1.5}
Let $X$ be a quasi-compact prescheme. If there exists an invertible
$\mathcal{O}_X$-module $\mathcal{L}$ such that $\mathcal{L}$ and
$\mathcal{L}^{-1}$ are ample, then $X$ is a quasi-affine scheme.
\end{corollary}

Indeed, $\mathcal{O}_X=\mathcal{L}\otimes\mathcal{L}^{-1}$ is then ample
(\hyperref[II.4.5.7]{4.5.7}).

\begin{proposition}[5.1.6]
\label{II.5.1.6}
Let $f:X\to Y$ be a quasi-compact morphism. The following conditions are
equivalent:
\begin{enumerate}
  \item[\rm{a)}] $f$ is quasi-affine.
  \item[\rm{b)}] The $\mathcal{O}_Y$-algebra
  $f_*(\mathcal{O}_X)=\mathcal{A}(X)$ is quasi-coherent and the canonical
  morphism $X\to\operatorname{Spec}(\mathcal{A}(X))$ corresponding to
  the identity homomorphism $\mathcal{A}(X)\to\mathcal{A}(X)$
  (\hyperref[II.1.2.7]{1.2.7}) is an open immersion.
  \item[\rm{b')}] The $\mathcal{O}_Y$-algebra $\mathcal{A}(X)$ is
  quasi-coherent and the canonical morphism
  $X\to\operatorname{Spec}(\mathcal{A}(X))$ is a homeomorphism of $X$
  onto a subspace of $\operatorname{Spec}(\mathcal{A}(X))$.
  \item[\rm{c)}] The $\mathcal{O}_X$-module $\mathcal{O}_X$ is very ample
  for $f$.
  \item[\rm{c')}] The $\mathcal{O}_X$-module $\mathcal{O}_X$ is ample for
  $f$.
  \item[\rm{d)}] The morphism $f$ is separated and, for every
  quasi-coherent $\mathcal{O}_X$-module $\mathcal{F}$, the canonical
  homomorphism $\sigma:f^*(f_*(\mathcal{F}))\to\mathcal{F}$
  (\hyperref[0.4.4.3]{0, 4.4.3}) is surjective.
\end{enumerate}
Moreover, when $f$ is quasi-affine, every invertible $\mathcal{O}_X$-module
$\mathcal{L}$ is very ample relative to $f$.
\end{proposition}

The equivalence of a) and c') follows from the local character on $Y$ of
$f$-ampleness (\hyperref[II.4.6.4]{4.6.4}), from Definition
(\hyperref[II.5.1.1]{5.1.1}) and from criterion
(\hyperref[II.5.1.2]{5.1.2, c')}). The other properties are local on
$Y$

\oldpage[II]{96}

and therefore follow immediately from (\hyperref[II.5.1.2]{5.1.2}) and
(\hyperref[II.5.1.4]{5.1.4}), taking into account the fact that
$f_*(\mathcal{F})$ is quasi-coherent when $f$ is separated
(\hyperref[I.9.2.2]{I, 9.2.2, a)}).

\begin{corollary}[5.1.7]
\label{II.5.1.7}
Let $f:X\to Y$ be a quasi-affine morphism. For every open subset $U$ of $Y$,
the restriction $f^{-1}(U)\to U$ of $f$ is quasi-affine.
\end{corollary}

\begin{corollary}[5.1.8]
\label{II.5.1.8}
Let $Y$ be an affine scheme and $f:X\to Y$ a quasi-compact morphism. For
$f$ to be quasi-affine, it is necessary and sufficient that $X$ be a
quasi-affine scheme.
\end{corollary}

This is an immediate consequence of (\hyperref[II.5.1.6]{5.1.6}) and
(\hyperref[II.4.6.6]{4.6.6}).

\begin{corollary}[5.1.9]
\label{II.5.1.9}
Let $Y$ be a quasi-compact scheme or a prescheme whose underlying space
is Noetherian, and $f:X\to Y$ a morphism of finite type. If $f$ is
quasi-affine, there exists a quasi-coherent sub-$\mathcal{O}_Y$-algebra
$\mathcal{B}$ of $\mathcal{A}(X)=f_*(\mathcal{O}_X)$, of finite type
(\hyperref[I.9.6.2]{I, 9.6.2}), such that the morphism
$X\to\operatorname{Spec}(\mathcal{B})$ corresponding to the canonical injection
$\mathcal{B}\to\mathcal{A}(X)$ (\hyperref[II.1.2.7]{1.2.7}) is an
immersion. Moreover, every quasi-coherent sub-$\mathcal{O}_Y$-algebra of
finite type $\mathcal{B}'$ of $\mathcal{A}(X)$ containing $\mathcal{B}$
has the same property.
\end{corollary}

Indeed, $\mathcal{A}(X)$ is the inductive limit of its quasi-coherent
sub-$\mathcal{O}_Y$-algebras of finite type
(\hyperref[I.9.6.5]{I, 9.6.5}); the result is then a particular case
of (\hyperref[II.3.8.4]{3.8.4}), taking into account the identification of
$\operatorname{Spec}(\mathcal{A}(X))$ with
$\operatorname{Proj}(\mathcal{A}(X)[T])$
(\hyperref[II.3.1.7]{3.1.7}).

\begin{proposition}[5.1.10]
\label{II.5.1.10}
\begin{enumerate}
  \item[\rm{(i)}] A quasi-compact morphism $X\to Y$ which is a
  homeomorphism of the underlying space of $X$ onto a subspace of the
  underlying space of $Y$ (in particular a closed immersion) is
  quasi-affine.
  \item[\rm{(ii)}] The composite of two quasi-affine morphisms is
  quasi-affine.
  \item[\rm{(iii)}] If $f:X\to Y$ is a quasi-affine $S$-morphism, then
  $f_{(S')}:X_{(S')}\to Y_{(S')}$ is a quasi-affine morphism for every
  extension $S'\to S$ of the base prescheme.
  \item[\rm{(iv)}] If $f:X\to Y$ and $g:X'\to Y'$ are two quasi-affine
  $S$-morphisms, then $f\times_S g$ is quasi-affine.
  \item[\rm{(v)}] If $f:X\to Y$, $g:Y\to Z$ are two morphisms such that
  $g\circ f$ is quasi-affine, and if $g$ is separated or the underlying
  space of $X$ is locally Noetherian, then $f$ is quasi-affine.
  \item[\rm{(vi)}] If $f$ is a quasi-affine morphism, then so is
  $f_{\mathrm{red}}$.
\end{enumerate}
\end{proposition}

Taking into account criterion (\hyperref[II.5.1.6]{5.1.6, c')}), (i),
(iii), (iv), (v) and (vi) follow immediately from
(\hyperref[II.4.6.13]{4.6.13, (i bis), (iii), (iv), (v) and (vi)})
respectively. To prove (ii), we may restrict to the case where $Z$ is
affine, and then the assertion follows directly from
(\hyperref[II.4.6.13]{4.6.13, (ii)}), applied to
$\mathcal{L}=\mathcal{O}_X$ and $\mathcal{K}=\mathcal{O}_Y$.

\begin{remark}[5.1.11]
\label{II.5.1.11}
Let $f:X\to Y$, $g:Y\to Z$ be two morphisms such that $X\times_ZY$ is
locally Noetherian. Then the graph immersion
$\Gamma_f:X\to X\times_ZY$ is quasi-affine, being quasi-compact
(\hyperref[I.6.3.5]{I, 6.3.5}), and
(\hyperref[I.5.5.12]{I, 5.5.12}) shows that, in (v), the conclusion remains
valid if we drop the hypothesis that $g$ is separated.
\end{remark}

\begin{proposition}[5.1.12]
\label{II.5.1.12}
Let $f:X\to Y$ be a quasi-compact morphism and $g:X'\to X$ a quasi-affine
morphism. If $\mathcal{L}$ is an $\mathcal{O}_X$-module ample for
$f$, then $g^*(\mathcal{L})$ is an $\mathcal{O}_{X'}$-module ample for
$f\circ g$.
\end{proposition}

Indeed, since $\mathcal{O}_{X'}$ is very ample for $g$, and the
question is local on $Y$ (\hyperref[II.4.6.4]{4.6.4}), it follows from
(\hyperref[II.4.6.13]{4.6.13, (ii)}) that there exists (for $Y$ affine) an
integer $n$ such that $g^*(\mathcal{L})^{\otimes n}$ is ample for
$f\circ g$, hence $g^*(\mathcal{L})$ is ample for $f\circ g$
(\hyperref[II.4.6.9]{4.6.9}).

\oldpage[II]{97}

\subsection{Serre's criterion.}
\label{subsection:II.5.2}

\begin{theorem}[5.2.1]
\label{II.5.2.1}
\emph{(Serre's criterion.)} Let $X$ be a quasi-compact scheme or a prescheme
whose underlying space is Noetherian. The following conditions are
equivalent:
\begin{enumerate}
  \item[\rm{a)}] $X$ is an affine scheme.
  \item[\rm{b)}] There exists a family of elements
  $f_\alpha\in A=\Gamma(X,\mathcal{O}_X)$ such that the $X_{f_\alpha}$ are
  affine and the ideal generated by the $f_\alpha$ in $A$ is equal to
  $A$.
  \item[\rm{c)}] The functor $\Gamma(X,\mathcal{F})$ is exact in
  $\mathcal{F}$ on the category of quasi-coherent
  $\mathcal{O}_X$-modules; in other words, if
  \[
    0\to\mathcal{F}'\to\mathcal{F}\to\mathcal{F}''\to0 \tag{*}
  \]
  is an exact sequence of quasi-coherent $\mathcal{O}_X$-modules, the
  sequence
  \[
    0\to\Gamma(X,\mathcal{F}')\to\Gamma(X,\mathcal{F})
      \to\Gamma(X,\mathcal{F}'')\to0
  \]
  is exact.
  \item[\rm{c')}] Property c) holds for every exact sequence (*) of
  quasi-coherent $\mathcal{O}_X$-modules such that $\mathcal{F}$ is
  isomorphic to a sub-$\mathcal{O}_X$-module of a finite product
  $\mathcal{O}_X^n$.
  \item[\rm{d)}] $\mathrm{H}^1(X,\mathcal{F})=0$ for every
  quasi-coherent $\mathcal{O}_X$-module $\mathcal{F}$.
  \item[\rm{d')}] $\mathrm{H}^1(X,\mathcal{J})=0$ for every
  quasi-coherent sheaf of ideals $\mathcal{J}$ of $\mathcal{O}_X$.
\end{enumerate}
\end{theorem}

It is evident that a) implies b); on the other hand, b) implies that the
$X_{f_\alpha}$ cover $X$, since by hypothesis the section $1$ is a
linear combination of the $f_\alpha$ and the $D(f_\alpha)$ then cover
$\operatorname{Spec}(A)$. The last assertion of
(\hyperref[II.4.5.2]{4.5.2}) therefore implies that
$X\to\operatorname{Spec}(A)$ is an isomorphism.

We know that a) implies c) (\hyperref[I.1.3.11]{I, 1.3.11}) and c)
trivially implies c'). Let us prove that c') implies b). First of all, c')
implies that, for every closed point $x\in X$ and every open neighbourhood
$U$ of $x$, there exists $f\in A$ such that $x\in X_f\subset X-U$. Indeed, let
$\mathcal{J}$ (resp. $\mathcal{J}'$) be the quasi-coherent sheaf of ideals of
$\mathcal{O}_X$ defining the reduced closed sub-prescheme of $X$ having
$X-U$ (resp. $(X-U)\cup\{x\}$) as underlying space
(\hyperref[I.5.2.1]{I, 5.2.1}); it is clear that we have
$\mathcal{J}'\subset\mathcal{J}$ and that
$\mathcal{J}''=\mathcal{J}/\mathcal{J}'$ is a quasi-coherent
$\mathcal{O}_X$-module having support $\{x\}$ and such that
$\mathcal{J}''_x=k(x)$. Hypothesis c') applied to the exact sequence
$0\to\mathcal{J}'\to\mathcal{J}\to\mathcal{J}''\to0$ shows that
$\Gamma(X,\mathcal{J})\to\Gamma(X,\mathcal{J}'')$ is surjective. The section
of $\mathcal{J}''$ whose germ at $x$ is $1_x$ is therefore the image of a
section $f\in\Gamma(X,\mathcal{J})\subset\Gamma(X,\mathcal{O}_X)$, and we have
by definition $f(x)=1_x$ and $f(y)=0$ in $X-U$, which establishes our
assertion. Moreover, if $U$ is affine, then so is $X_f$
(\hyperref[I.1.3.6]{I, 1.3.6}), hence the union of those $X_f$ which are
affine ($f\in A$) is an open set $Z$ containing all the closed points of
$X$; since $X$ is a quasi-compact Kolmogorov space, we necessarily have
$Z=X$ (\hyperref[0.2.1.3]{0, 2.1.3}). Since $X$ is
quasi-compact, there are finitely many elements $f_i\in A$
($1\leq i\leq n$) such that the $X_{f_i}$ are affine and cover $X$.
Consider then the homomorphism $\mathcal{O}_X^n\to\mathcal{O}_X$ defined
by the sections $f_i$ (\hyperref[0.5.1.1]{0, 5.1.1}); since, for every
$x\in X$, at least one of the $(f_i)_x$ is invertible, this homomorphism is
surjective, and we therefore have an exact sequence
$0\to\mathcal{R}\to\mathcal{O}_X^n\to\mathcal{O}_X\to0$, where $\mathcal{R}$
is a quasi-coherent sub-$\mathcal{O}_X$-module of $\mathcal{O}_X^n$. It
then follows

\oldpage[II]{98}

from c') that the corresponding homomorphism
$\Gamma(X,\mathcal{O}_X^n)\to\Gamma(X,\mathcal{O}_X)$ is surjective, which
proves b).

Finally, a) implies d) (\hyperref[I.5.1.9.2]{I, 5.1.9.2}) and d)
trivially implies d'). It remains to show that d') implies c'). Now, if
$\mathcal{F}'$ is a quasi-coherent sub-$\mathcal{O}_X$-module of
$\mathcal{O}_X^n$, the filtration
$0\subset\mathcal{O}_X\subset\mathcal{O}_X^2\subset\cdots\subset
\mathcal{O}_X^n$ defines on $\mathcal{F}'$ a filtration formed by the
$\mathcal{F}'_k=\mathcal{O}_X^k\cap\mathcal{F}'$ ($0\leq k\leq n$), which
are quasi-coherent $\mathcal{O}_X$-modules
(\hyperref[I.4.1.1]{I, 4.1.1}), and
$\mathcal{F}'_{k+1}/\mathcal{F}'_k$ is isomorphic to a quasi-coherent
sub-$\mathcal{O}_X$-module of
$\mathcal{O}_X^{k+1}/\mathcal{O}_X^k=\mathcal{O}_X$, that is, to a
quasi-coherent sheaf of ideals of $\mathcal{O}_X$. Hypothesis d')
therefore implies
$\mathrm{H}^1(X,\mathcal{F}'_{k+1}/\mathcal{F}'_k)=0$; the exact cohomology
sequence
$\mathrm{H}^1(X,\mathcal{F}'_k)\to\mathrm{H}^1(X,\mathcal{F}'_{k+1})
\to\mathrm{H}^1(X,\mathcal{F}'_{k+1}/\mathcal{F}'_k)=0$ then allows us to
prove by induction on $k$ that
$\mathrm{H}^1(X,\mathcal{F}'_k)=0$ for every $k$. Q.E.D.

\begin{remark}[5.2.1.1]
\label{II.5.2.1.1}
When $X$ is a Noetherian prescheme, we may, in the statement of c') and of
d'), replace ``quasi-coherent'' by ``coherent''. Indeed, in the proof of
the fact that c') implies b), $\mathcal{J}$ and $\mathcal{J}'$
are then coherent sheaves of ideals, and, on the other hand, every
quasi-coherent submodule of a coherent module is coherent
(\hyperref[I.6.1.1]{I, 6.1.1}); whence the conclusion.
\end{remark}

\begin{corollary}[5.2.2]
\label{II.5.2.2}
Let $f:X\to Y$ be a separated quasi-compact morphism. The following
conditions are equivalent:
\begin{enumerate}
  \item[\rm{a)}] $f$ is an affine morphism.
  \item[\rm{b)}] The functor $f_*$ is exact on the category of
  quasi-coherent $\mathcal{O}_X$-modules.
  \item[\rm{c)}] For every quasi-coherent $\mathcal{O}_X$-module
  $\mathcal{F}$, we have $\mathrm{R}^1f_*(\mathcal{F})=0$.
  \item[\rm{c')}] For every quasi-coherent sheaf of ideals $\mathcal{J}$
  of $\mathcal{O}_X$, we have $\mathrm{R}^1f_*(\mathcal{J})=0$.
\end{enumerate}
\end{corollary}

All these conditions being local on $Y$, by definition of the functor
$\mathrm{R}^1f_*$ (T, 3.7.3), we may suppose that $Y$ is affine. If $f$
is affine, $X$ is then affine and property b) is none other than
(\hyperref[I.1.6.4]{I, 1.6.4}). Conversely, let us show that b) implies
a): for every quasi-coherent $\mathcal{O}_X$-module $\mathcal{F}$,
$f_*(\mathcal{F})$ is a quasi-coherent $\mathcal{O}_Y$-module
(\hyperref[I.9.2.2]{I, 9.2.2, a)}). By hypothesis the functor
$f_*(\mathcal{F})$ is exact in $\mathcal{F}$, and the functor
$\Gamma(Y,\mathcal{G})$ is exact in $\mathcal{G}$ (on the category of
quasi-coherent $\mathcal{O}_Y$-modules) since $Y$ is affine
(\hyperref[I.1.3.11]{I, 1.3.11}); hence
$\Gamma(Y,f_*(\mathcal{F}))=\Gamma(X,\mathcal{F})$ is exact in
$\mathcal{F}$, which proves our assertion by virtue of
(\hyperref[II.5.2.1]{5.2.1, c)}).

If $f$ is affine, $f^{-1}(U)$ is affine for every affine open subset $U$ of
$Y$ (\hyperref[II.1.3.2]{1.3.2}), hence
$\mathrm{H}^1(f^{-1}(U),\mathcal{F})=0$
(\hyperref[II.5.2.1]{5.2.1, d)}), which by definition implies
$\mathrm{R}^1f_*(\mathcal{F})=0$. Finally, suppose that condition
c') is satisfied; the exact sequence of terms of low degree in the
Leray spectral sequence (G, II, 4.17.1 and I, 4.5.1) gives in particular
the exact sequence
\[
  0\to\mathrm{H}^1(Y,f_*(\mathcal{J}))\to
  \mathrm{H}^1(X,\mathcal{J})\to
  \mathrm{H}^0(Y,\mathrm{R}^1f_*(\mathcal{J})).
\]
Since $Y$ is affine and $f_*(\mathcal{J})$ is quasi-coherent
(\hyperref[I.9.2.2]{I, 9.2.2, a)}), we have
$\mathrm{H}^1(Y,f_*(\mathcal{J}))=0$
(\hyperref[II.5.2.1]{5.2.1}); hypothesis c') therefore implies
$\mathrm{H}^1(X,\mathcal{J})=0$, and we conclude by
(\hyperref[II.5.2.1]{5.2.1}) that $X$ is an affine scheme.

\begin{corollary}[5.2.3]
\label{II.5.2.3}
If $f:X\to Y$ is an affine morphism, then, for every quasi-coherent
$\mathcal{O}_X$-module $\mathcal{F}$, the canonical homomorphism
$\mathrm{H}^1(Y,f_*(\mathcal{F}))\to\mathrm{H}^1(X,\mathcal{F})$ is
bijective.
\end{corollary}

\oldpage[II]{99}

Indeed, we have the exact sequence
\[
  0\to\mathrm{H}^1(Y,f_*(\mathcal{F}))\to
  \mathrm{H}^1(X,\mathcal{F})\to
  \mathrm{H}^0(Y,\mathrm{R}^1f_*(\mathcal{F}))
\]
of terms of low degree of the Leray spectral sequence, and the conclusion
follows from (\hyperref[II.5.2.2]{5.2.2}).

\begin{remark}[5.2.4]
\label{II.5.2.4}
In Chapter III, § 1, we will prove that if $X$ is affine, we have
$\mathrm{H}^i(X,\mathcal{F})=0$ for every $i>0$ and every quasi-coherent
$\mathcal{O}_X$-module $\mathcal{F}$.
\end{remark}

\subsection{Quasi-projective morphisms.}
\label{subsection:II.5.3}

\begin{definition}[5.3.1]
\label{II.5.3.1}
We say that a morphism $f:X\to Y$ is \emph{quasi-projective}, or that $X$
(considered as a $Y$-prescheme by means of $f$) is \emph{quasi-projective
over $Y$}, or that $X$ is a \emph{quasi-projective $Y$-scheme}, if $f$ is of
finite type and there exists an invertible $\mathcal{O}_X$-module which is
$f$-ample.
\end{definition}

We note that this notion is not local on $Y$: the counterexamples of
Nagata [26] and of Hironaka show that, even if $X$ and $Y$ are nonsingular
algebraic schemes over an algebraically closed field, every point of
$Y$ may have an affine neighbourhood $U$ such that $f^{-1}(U)$ is
quasi-projective over $U$, without $f$ being quasi-projective.

We note that a quasi-projective morphism is necessarily separated
(\hyperref[II.4.6.1]{4.6.1}). When $Y$ is quasi-compact, it amounts to
the same thing to say that $f$ is quasi-projective or that $f$ is of finite
type and there exists an $\mathcal{O}_X$-module very ample relative to $f$
(\hyperref[II.4.6.2]{4.6.2} and \hyperref[II.4.6.11]{4.6.11}). Moreover:

\begin{proposition}[5.3.2]
\label{II.5.3.2}
Let $Y$ be a quasi-compact scheme or a prescheme whose underlying space
is Noetherian, and let $X$ be a $Y$-prescheme. The following conditions are
equivalent:
\begin{enumerate}
  \item[\rm{a)}] $X$ is a quasi-projective $Y$-scheme.
  \item[\rm{b)}] $X$ is of finite type over $Y$, and there exists a
  quasi-coherent $\mathcal{O}_Y$-module of finite type $\mathcal{E}$ such that
  $X$ is $Y$-isomorphic to a sub-prescheme of $\mathbf{P}(\mathcal{E})$.
  \item[\rm{c)}] $X$ is of finite type over $Y$, and there exists a
  quasi-coherent graded $\mathcal{O}_Y$-algebra $\mathcal{S}$ such that
  $\mathcal{S}_1$ is of finite type and generates $\mathcal{S}$ and such that
  $X$ is $Y$-isomorphic to a sub-prescheme induced on an everywhere dense
  open subset of $\operatorname{Proj}(\mathcal{S})$.
\end{enumerate}
\end{proposition}

This follows immediately from the preceding remark and from
(\hyperref[II.4.4.3]{4.4.3}), (\hyperref[II.4.4.6]{4.4.6}) and
(\hyperref[II.4.4.7]{4.4.7}).

We note that when $Y$ is a Noetherian prescheme, we may, in
conditions b) and c) of (\hyperref[II.5.3.2]{5.3.2}), drop the
hypothesis that $X$ is of finite type over $Y$, which is then automatically
satisfied (\hyperref[I.6.3.5]{I, 6.3.5}).

\begin{corollary}[5.3.3]
\label{II.5.3.3}
Let $Y$ be a quasi-compact scheme such that there exists an ample
$\mathcal{O}_Y$-module $\mathcal{L}$ (\hyperref[II.4.5.3]{4.5.3}). For a
$Y$-scheme $X$ to be quasi-projective, it is necessary and sufficient that it
be of finite type over $Y$ and isomorphic to a sub-$Y$-scheme of a projective
bundle of the form $\mathbf{P}_Y^r$.
\end{corollary}

Indeed, if $\mathcal{E}$ is a quasi-coherent $\mathcal{O}_Y$-module of
finite type, $\mathcal{E}$ is isomorphic to a quotient of an
$\mathcal{O}_Y$-module
$\mathcal{L}^{\otimes(-n)}\otimes_{\mathcal{O}_Y}\mathcal{O}_Y^k$
(\hyperref[II.4.5.5]{4.5.5}), hence $\mathbf{P}(\mathcal{E})$ is
isomorphic to a closed subscheme of $\mathbf{P}_Y^{k-1}$
(\hyperref[II.4.1.2]{4.1.2} and \hyperref[II.4.1.4]{4.1.4}).

\begin{proposition}[5.3.4]
\label{II.5.3.4}
\begin{enumerate}
  \item[\rm{(i)}] Every quasi-affine morphism of finite type (in particular every
  quasi-compact immersion and every affine morphism of finite type) is
  quasi-projective.
  \item[\rm{(ii)}] If $f:X\to Y$ and $g:Y\to Z$ are quasi-projective and if
  $Z$ is quasi-compact, then $g\circ f$ is quasi-projective.

\oldpage[II]{100}

  \item[\rm{(iii)}] If $f:X\to Y$ is a quasi-projective $S$-morphism, then
  $f_{(S')}:X_{(S')}\to Y_{(S')}$ is quasi-projective for every extension
  $S'\to S$ of the base prescheme.
  \item[\rm{(iv)}] If $f:X\to Y$ and $g:X'\to Y'$ are two quasi-projective
  $S$-morphisms, then $f\times_Sg$ is quasi-projective.
  \item[\rm{(v)}] If $f:X\to Y$, $g:Y\to Z$ are two morphisms such that
  $g\circ f$ is quasi-projective, and if $g$ is separated or $X$ is locally
  Noetherian, then $f$ is quasi-projective.
  \item[\rm{(vi)}] If $f$ is a quasi-projective morphism, then so is
  $f_{\mathrm{red}}$.
\end{enumerate}
\end{proposition}

(i) follows from (\hyperref[II.5.1.6]{5.1.6}) and
(\hyperref[II.5.1.10]{5.1.10, (i)}). The other assertions are immediate
consequences of Definition (\hyperref[II.5.3.1]{5.3.1}), of the properties
of morphisms of finite type (\hyperref[I.6.3.4]{I, 6.3.4}) and of
(\hyperref[II.4.6.13]{4.6.13}).

\begin{remark}[5.3.5]
\label{II.5.3.5}
We note that it can happen that $f_{\mathrm{red}}$ is quasi-projective
without $f$ being so, even assuming that $Y$ is the spectrum of an algebra
of finite rank over $\mathbf{C}$ and that $f$ is proper.
\end{remark}

\begin{corollary}[5.3.6]
\label{II.5.3.6}
If $X$ and $X'$ are two quasi-projective $Y$-schemes, then $X\amalg X'$ is a
quasi-projective $Y$-scheme.
\end{corollary}

This follows from (\hyperref[II.4.6.18]{4.6.18}).

\subsection{Proper morphisms and universally closed morphisms.}
\label{subsection:II.5.4}
\label{II.5.4}

\begin{definition}[5.4.1]
\label{II.5.4.1}
We say that a morphism of preschemes $f:X\to Y$ is \emph{proper} if it
satisfies the following two conditions:
\begin{enumerate}
  \item[\rm{a)}] $f$ is separated and of finite type.
  \item[\rm{b)}] For every prescheme $Y'$ and every morphism $Y'\to Y$, the
  projection $f_{(Y')}:X\times_YY'\to Y'$ is a closed morphism
  (\hyperref[I.2.2.6]{I, 2.2.6}).
\end{enumerate}
When this is the case, we also say that $X$ (considered as a $Y$-prescheme
with structural morphism $f$) is proper over $Y$.
\end{definition}

It is immediate that conditions a) and b) are local on $Y$. To verify
that the image of a closed subset $Z$ of $X\times_YY'$ under the
projection $q:X\times_YY'\to Y'$ is closed in $Y'$, it suffices to see that
$q(Z)\cap U'$ is closed in $U'$ for every affine open subset $U'$ of $Y'$;
since $q(Z)\cap U'=q(Z\cap q^{-1}(U'))$ and $q^{-1}(U')$ is identified with
$X\times_YU'$ (\hyperref[I.4.4.1]{I, 4.4.1}), we see that to verify
condition b) of Def. (\hyperref[II.5.4.1]{5.4.1}), we may restrict to the
case where $Y'$ is an affine scheme. We will see later
(\hyperref[II.5.6.3]{5.6.3}) that if $Y$ is locally Noetherian, we
may even restrict ourselves to verifying b) when $Y'$ is of finite type over $Y$.

It is clear that every proper morphism is \emph{closed}.

\begin{proposition}[5.4.2]
\label{II.5.4.2}
\begin{enumerate}
  \item[\rm{(i)}] A closed immersion is a proper morphism.
  \item[\rm{(ii)}] The composite of two proper morphisms is proper.
  \item[\rm{(iii)}] If $X$, $Y$ are $S$-preschemes and
  $f:X\to Y$ is a proper $S$-morphism, then
  $f_{(S')}:X_{(S')}\to Y_{(S')}$ is proper for every extension
  $S'\to S$ of the base prescheme.
  \item[\rm{(iv)}] If $f:X\to Y$ and $g:X'\to Y'$ are two proper
  $S$-morphisms, the $S$-morphism
  $f\times_Sg:X\times_SX'\to Y\times_SY'$ is proper.
\end{enumerate}
\end{proposition}

\oldpage[II]{101}

It suffices to prove (i), (ii) and (iii) (\hyperref[I.3.5.1]{I, 3.5.1}).
In each of these three cases, the verification of condition a) of
(\hyperref[II.5.4.1]{5.4.1}) follows from earlier results
(\hyperref[I.5.5.1]{I, 5.5.1} and \hyperref[I.6.3.4]{6.3.4}); it remains to
verify condition b). This is immediate in case (i), for if $X\to Y$
is a closed immersion, then so is
$X\times_YY'\to Y\times_YY'=Y'$ (\hyperref[I.4.3.2]{I, 4.3.2} and
\hyperref[II.3.3.3]{3.3.3}). To prove (ii), consider two
proper morphisms $X\to Y$, $Y\to Z$, and a morphism $Z'\to Z$. We can
write $X\times_ZZ'=X\times_Y(Y\times_ZZ')$
(\hyperref[I.3.3.9.1]{I, 3.3.9.1}), and consequently the projection
$X\times_ZZ'\to Z'$ factors as
$X\times_Y(Y\times_ZZ')\to Y\times_ZZ'\to Z'$. Taking into account the remark
at the beginning, (ii) therefore follows from the fact that the composite of
two closed morphisms is closed. Finally, for every morphism $S'\to S$,
$X_{(S')}$ is identified with
$X\times_YY_{(S')}$ (\hyperref[I.3.3.11]{I, 3.3.11}); for every
morphism $Z\to Y_{(S')}$, we can write
\[
  X_{(S')}\times_{Y_{(S')}}Z
  =(X\times_YY_{(S')})\times_{Y_{(S')}}Z=X\times_YZ~;
\]
$X\times_YZ\to Z$ is closed by hypothesis, so this proves (iii).

\begin{corollary}[5.4.3]
\label{II.5.4.3}
Let $f:X\to Y$, $g:Y\to Z$ be two morphisms such that $g\circ f$ is
proper.
\begin{enumerate}
  \item[\rm{(i)}] If $g$ is separated, then $f$ is proper.
  \item[\rm{(ii)}] If $g$ is separated and of finite type, and if $f$ is
  surjective, then $g$ is proper.
\end{enumerate}
\end{corollary}

(i) follows from (\hyperref[II.5.4.2]{5.4.2}) by the general procedure
(\hyperref[I.5.5.12]{I, 5.5.12}). To prove (ii), we need only
verify condition b) of Def. (\hyperref[II.5.4.1]{5.4.1}).
For every morphism $Z'\to Z$, the diagram
\[
  \xymatrix{
    X\times_ZZ'\ar[r]^{f\times1_{Z'}}\ar[dr]_p &
    Y\times_ZZ'\ar[d]^{p'}\\
    & Z'
  }
\]
(where $p$ and $p'$ are the projections) is commutative
(\hyperref[I.3.2.1]{I, 3.2.1}); moreover, $f\times1_{Z'}$ is surjective
since $f$ is (\hyperref[I.3.5.2]{I, 3.5.2}), and $p$ is
a closed morphism by hypothesis. Every closed subset $F$ of
$Y\times_ZZ'$ is then the image under $f\times1_{Z'}$ of a closed subset $E$
of $X\times_ZZ'$, hence $p'(F)=p(E)$ is closed in $Z'$ by hypothesis, whence
the corollary.

\begin{corollary}[5.4.4]
\label{II.5.4.4}
If $X$ is a prescheme proper over $Y$, and $\mathcal{S}$ a
quasi-coherent $\mathcal{O}_Y$-algebra, every $Y$-morphism
$f:X\to\operatorname{Proj}(\mathcal{S})$ is proper (and \emph{a fortiori}
closed).
\end{corollary}

Indeed, the structural morphism
$p:\operatorname{Proj}(\mathcal{S})\to Y$ is separated, and $p\circ f$ is
proper by hypothesis.

\begin{corollary}[5.4.5]
\label{II.5.4.5}
Let $f:X\to Y$ be a separated morphism of finite type. Let
$(X_i)_{1\leq i\leq n}$ (resp. $(Y_i)_{1\leq i\leq n}$) be a finite family
of closed sub-preschemes of $X$ (resp. $Y$), and $j_i$ (resp. $h_i$)
the canonical injection $X_i\to X$ (resp. $Y_i\to Y$). Suppose that the
underlying space $X$ is the union of the $X_i$, and that for every $i$ there
is a morphism $f_i:X_i\to Y_i$ such that the diagram
\[
  \xymatrix{
    X_i\ar[r]^{f_i}\ar[d]_{j_i} & Y_i\ar[d]^{h_i}\\
    X\ar[r]_f & Y
  }
\]
is commutative. Then, for $f$ to be proper, it is necessary and sufficient
that each of the $f_i$ be proper.
\end{corollary}

\oldpage[II]{102}

If $f$ is proper, then so is $f\circ j_i$, since $j_i$ is a
closed immersion (\hyperref[II.5.4.2]{5.4.2}); since $h_i$ is a
closed immersion, hence a separated morphism, $f_i$ is proper by
(\hyperref[II.5.4.3]{5.4.3}). Conversely, suppose that each of the $f_i$
is proper, and consider the prescheme $Z$, the sum of the $X_i$; let $u$ be
the morphism $Z\to X$ which reduces to $j_i$ on each of the $X_i$. The
restriction of $f\circ u$ to each $X_i$ is equal to
$f\circ j_i=h_i\circ f_i$, hence is proper since $h_i$
and $f_i$ are (\hyperref[II.5.4.2]{5.4.2}); it then follows immediately
from Def. (\hyperref[II.5.4.1]{5.4.1}) that $f\circ u$ is proper. But
since $u$ is surjective by hypothesis, we conclude that $f$ is proper by
(\hyperref[II.5.4.3]{5.4.3}).

\begin{corollary}[5.4.6]
\label{II.5.4.6}
Let $f:X\to Y$ be a separated morphism of finite type; for $f$ to be
proper, it is necessary and sufficient that
$f_{\mathrm{red}}:X_{\mathrm{red}}\to Y_{\mathrm{red}}$ be proper.
\end{corollary}

This is a particular case of (\hyperref[II.5.4.5]{5.4.5}) with $n=1$,
$X_1=X_{\mathrm{red}}$ and $Y_1=Y_{\mathrm{red}}$
(\hyperref[I.5.1.5]{I, 5.1.5}).

\begin{env}[5.4.7]
\label{II.5.4.7}
If $X$ and $Y$ are Noetherian preschemes and $f:X\to Y$ is a
separated morphism of finite type, we can, in order to verify that $f$ is
proper, reduce to doing so for dominant morphisms of integral preschemes.
Indeed, let $X_i$ ($1\leq i\leq n$) be the irreducible components of $X$,
which are finite in number, and consider for each $i$ the unique reduced
closed sub-prescheme of $X$ having $X_i$ as underlying space, which we will
again denote by $X_i$ (\hyperref[I.5.2.1]{I, 5.2.1}). On the other hand,
let $Y_i$ be the unique reduced closed sub-prescheme of $Y$ having
$\overline{f(X_i)}$ as underlying space. If $g_i$ (resp. $h_i$) is the
injection morphism $X_i\to X$ (resp. $Y_i\to Y$), we conclude that
$f\circ g_i=h_i\circ f_i$, where $f_i$ is a dominant morphism $X_i\to Y_i$
(\hyperref[I.5.2.2]{I, 5.2.2}); we are then in a position to apply
(\hyperref[II.5.4.5]{5.4.5}), and for $f$ to be proper, it is necessary
and sufficient that the $f_i$ be proper.
\end{env}

\begin{corollary}[5.4.8]
\label{II.5.4.8}
Let $X$, $Y$ be $S$-preschemes that are separated and of finite type over
$S$, and let $f:X\to Y$ be an $S$-morphism. For $f$ to be proper, it is
necessary and sufficient that, for every $S$-prescheme $S'$, the morphism
$f\times_S1_{S'}:X\times_SS'\to Y\times_SS'$ be closed.
\end{corollary}

Let us note first that if $g:X\to S$ and $h:Y\to S$ are the structural
morphisms, we have by definition $g=h\circ f$, hence $f$ is separated and of
finite type (\hyperref[I.5.5.1]{I, 5.5.1} and
\hyperref[I.6.3.4]{6.3.4}). If $f$ is proper, then so is
$f\times_S1_{S'}$ (\hyperref[II.5.4.2]{5.4.2}); \emph{a fortiori},
$f\times_S1_{S'}$ is closed. Conversely, suppose that the condition of the
statement is satisfied and let $Y'$ be a $Y$-prescheme; $Y'$ can also be
considered as an $S$-prescheme, and since $Y\to S$ is separated,
$X\times_YY'$ is identified with a closed sub-prescheme of $X\times_SY'$
(\hyperref[I.5.4.2]{I, 5.4.2}). In the commutative diagram
\[
  \xymatrix{
    X\times_YY'\ar[r]^{f\times1_{Y'}}\ar[d] &
    Y\times_YY'=Y'\ar[d]\\
    X\times_SY'\ar[r]_{f\times_S1_{Y'}} & Y\times_SY'
  }
\]
the vertical arrows are closed immersions; it follows immediately
that if $f\times_S1_{Y'}$ is a closed morphism, then so is
$f\times_Y1_{Y'}$.

\begin{remark}[5.4.9]
\label{II.5.4.9}
We will say that a morphism $f:X\to Y$ is \emph{universally closed} if it
satisfies condition b) of Definition
(\hyperref[II.5.4.1]{5.4.1}). The reader will observe that

\oldpage[II]{103}

in Propositions (\hyperref[II.5.4.2]{5.4.2}) to
(\hyperref[II.5.4.8]{5.4.8}), we may everywhere replace ``proper'' by
``universally closed'' without affecting the validity of the results (and in
the hypotheses of (\hyperref[II.5.4.3]{5.4.3}),
(\hyperref[II.5.4.5]{5.4.5}), (\hyperref[II.5.4.6]{5.4.6}) and
(\hyperref[II.5.4.8]{5.4.8}), we may omit the finiteness
conditions).
\end{remark}

\begin{env}[5.4.10]
\label{II.5.4.10}
Let $f:X\to Y$ be a morphism of finite type. We will say that a closed subset
$Z$ of $X$ is \emph{proper over $Y$} (or \emph{$Y$-proper}, or \emph{proper
for $f$}) if the restriction of $f$ to a closed sub-prescheme of $X$
with underlying space $Z$ (\hyperref[I.5.2.1]{I, 5.2.1}) is proper.
Since this restriction is then separated, it follows from
(\hyperref[II.5.4.6]{5.4.6}) and (\hyperref[I.5.5.1]{I, 5.5.1, (vi)})
that the preceding property does not depend on the closed sub-prescheme of
$X$ having $Z$ as underlying space. If $g:X'\to X$ is a proper morphism,
$g^{-1}(Z)$ is a proper subset of $X'$: indeed, if $T$ is a sub-prescheme of
$X$ having $Z$ as underlying space, it suffices to note that the
restriction of $g$ to the closed sub-prescheme $g^{-1}(T)$ of $X'$ is a
proper morphism $g^{-1}(T)\to T$ by
(\hyperref[II.5.4.2]{5.4.2, (iii)}), and then to apply
(\hyperref[II.5.4.2]{5.4.2, (ii)}). On the other hand, if $X''$ is a
$Y$-scheme of finite type and $u:X\to X''$ a $Y$-morphism, $u(Z)$ is a
proper subset of $X''$; indeed, take for $T$ the reduced closed
sub-prescheme of $X$ having $Z$ as underlying space; since the restriction
of $f$ to $T$ is proper, so is the restriction of $u$ to $T$
(\hyperref[II.5.4.3]{5.4.3, (i)}), hence $u(Z)$ is closed in $X''$;
let $T''$ be a closed sub-prescheme of $X''$ having $u(Z)$ as underlying
space (\hyperref[I.5.2.1]{I, 5.2.1}), so that $u|T$ factors as
\[
  T\xrightarrow{v}T''\xrightarrow{j}X'',
\]
where $j$ is the canonical injection (\hyperref[I.5.2.2]{I, 5.2.2}) and $v$
is therefore proper and surjective (\hyperref[II.5.4.5]{5.4.5}); if $g$ is
the restriction to $T''$ of the structural morphism $X''\to Y$, then $g$ is
separated and of finite type, and we have $f|T=g\circ v$; it therefore
follows from (\hyperref[II.5.4.3]{5.4.3, (ii)}) that $g$ is proper, whence
our assertion.
\end{env}

It follows in particular from these remarks that if $Z$ is a
$Y$-proper subset of $X$, then:
\begin{enumerate}
  \item[\rm{1\textsuperscript{o}}] For every closed sub-prescheme $X'$ of
  $X$, $Z\cap X'$ is a $Y$-proper subset of $X'$.
  \item[\rm{2\textsuperscript{o}}] If $X$ is a sub-prescheme of a
  $Y$-scheme of finite type $X''$, then $Z$ is also a $Y$-proper subset of
  $X''$ (and in particular is closed in $X''$).
\end{enumerate}


\subsection{Projective morphisms.}
\label{subsection:II.5.5}

\begin{proposition}[5.5.1]
\label{II.5.5.1}
Let $X$ be a $Y$-prescheme. The following conditions are equivalent:
\begin{enumerate}
  \item[\rm{a)}] $X$ is $Y$-isomorphic to a closed sub-prescheme of a
  projective bundle $\mathbf{P}(\mathcal{E})$, where $\mathcal{E}$ is a
  quasi-coherent $\mathcal{O}_Y$-module of finite type.
  \item[\rm{b)}] There exists a quasi-coherent graded $\mathcal{O}_Y$-algebra
  $\mathcal{S}$ such that $\mathcal{S}_1$ is of finite type
  and generates $\mathcal{S}$, and such that $X$ is $Y$-isomorphic to
  $\operatorname{Proj}(\mathcal{S})$.
\end{enumerate}
\end{proposition}

Condition a) implies b) by virtue of
(\hyperref[II.3.6.2]{3.6.2, (ii)}): if $\mathcal{J}$ is a quasi-coherent
graded sheaf of ideals of $\mathbf{S}(\mathcal{E})$, the
quasi-coherent graded $\mathcal{O}_Y$-algebra
$\mathcal{S}=\mathbf{S}(\mathcal{E})/\mathcal{J}$ is generated by
$\mathcal{S}_1$, and the latter, the canonical image of $\mathcal{E}$, is an
$\mathcal{O}_Y$-module of finite type. Condition b) implies a) by virtue of
(\hyperref[II.3.6.2]{3.6.2}) applied to the case where
$\mathcal{M}\to\mathcal{S}_1$ is the identity map.

\oldpage[II]{104}

\begin{definition}[5.5.2]
\label{II.5.5.2}
We say that a $Y$-prescheme $X$ is \emph{projective over $Y$}, or is a
\emph{projective $Y$-scheme}, if it satisfies the equivalent conditions a)
and b) of (\hyperref[II.5.5.1]{5.5.1}). We say that a morphism $f:X\to Y$
is \emph{projective} if it makes $X$ a projective $Y$-scheme.
\end{definition}

It is clear that if $f:X\to Y$ is projective, there exists an
$\mathcal{O}_X$-module \emph{very ample relative to $f$}
(\hyperref[II.4.4.2]{4.4.2}).

\begin{theorem}[5.5.3]
\label{II.5.5.3}
\begin{enumerate}
  \item[\rm{(i)}] Every projective morphism is quasi-projective and proper.
  \item[\rm{(ii)}] Conversely, let $Y$ be a quasi-compact scheme or a
  prescheme whose underlying space is Noetherian; then every morphism
  $f:X\to Y$ which is quasi-projective and proper is projective.
\end{enumerate}
\end{theorem}

(i) It is clear that if $f:X\to Y$ is projective, it is of finite type and
quasi-projective (hence in particular separated); on the other hand, it
follows immediately from (\hyperref[II.5.5.1]{5.5.1, b)}) and from
(\hyperref[II.3.5.3]{3.5.3}) that if $f$ is projective, then so is
$f\times_Y1_{Y'}:X\times_YY'\to Y'$ for every morphism $Y'\to Y$. To show
that $f$ is universally closed, everything therefore comes down to proving
that a projective morphism $f$ is \emph{closed}. The question being local on
$Y$, we may suppose that $Y=\operatorname{Spec}(A)$, hence
(\hyperref[II.5.5.1]{5.5.1}) $X=\operatorname{Proj}(S)$, where $S$ is a
graded $A$-algebra generated by finitely many elements of $S_1$. For
every $y\in Y$, the fibre $f^{-1}(y)$ is identified with
$\operatorname{Proj}(S)\times_Y\operatorname{Spec}(k(y))$
(\hyperref[I.3.6.1]{I, 3.6.1}), hence with
$\operatorname{Proj}(S\otimes_Ak(y))$
(\hyperref[II.2.8.10]{2.8.10}); consequently, $f^{-1}(y)$ is empty if and
only if $S\otimes_Ak(y)$ satisfies condition (TN)
(\hyperref[II.2.7.4]{2.7.4}), in other words $S_n\otimes_Ak(y)=0$ for $n$
sufficiently large. But since $(S_n)_y$ is an $\mathcal{O}_y$-module of
finite type, the preceding condition means that $(S_n)_y=0$ for $n$
sufficiently large, by virtue of Nakayama's lemma. If $\mathfrak a_n$ is the
annihilator in $A$ of the $A$-module $S_n$, the preceding condition also
means that $\mathfrak a_n\not\subset\mathfrak j_y$ for $n$ sufficiently large
(\hyperref[0.1.7.4]{0, 1.7.4}). Now, since $S_nS_1=S_{n+1}$ by hypothesis,
we have $\mathfrak a_n\subset\mathfrak a_{n+1}$, and if $\mathfrak a$ is the
union of the $\mathfrak a_n$, we therefore see that $f(X)=V(\mathfrak a)$,
which proves that $f(X)$ is closed in $Y$. If now $X'$ is an arbitrary closed
subset of $X$, there exists a closed sub-prescheme of $X$ having $X'$
as underlying space (\hyperref[I.5.2.1]{I, 5.2.1}) and it is clear
(\hyperref[II.5.5.1]{5.5.1, a)}) that the morphism
$X'\to X\xrightarrow{f}Y$ is projective, hence $f(X')$ is closed in $Y$.

(ii) The hypothesis on $Y$ and the fact that $f$ is quasi-projective imply
the existence of a quasi-coherent $\mathcal{O}_Y$-module of finite type
$\mathcal{E}$ and of a $Y$-immersion $j:X\to\mathbf{P}(\mathcal{E})$
(\hyperref[II.5.3.2]{5.3.2}). But since $f$ is proper, $j$ is closed by
(\hyperref[II.5.4.4]{5.4.4}), hence $f$ is projective.

\par\medskip\noindent\textit{Remarks (5.5.4). ---\ }
\egalabel{II.5.5.4}{5.5.4}
\begin{enumerate}
  \item[\rm{(i)}] Let $f:X\to Y$ be a morphism such that:
  1\textsuperscript{o} $f$ is proper; 2\textsuperscript{o} there exists an
  $\mathcal{O}_X$-module $\mathcal{L}$ very ample relative to $f$;
  3\textsuperscript{o} the quasi-coherent $\mathcal{O}_Y$-module
  $\mathcal{E}=f_*(\mathcal{L})$ is of finite type. Then $f$ is a projective
  morphism: indeed (\hyperref[II.4.4.4]{4.4.4}), there is then a
  $Y$-immersion $r:X\to\mathbf{P}(\mathcal{E})$, and since $f$ is proper, $r$
  is a closed immersion (\hyperref[II.5.4.4]{5.4.4}). We will see in
  Chapter III, \S~3, that when $Y$ is locally Noetherian, condition
  3\textsuperscript{o} above is a consequence of the other two, so that
  conditions 1\textsuperscript{o} and 2\textsuperscript{o} characterize
  projective morphisms in this case, and if $Y$ is quasi-compact, we may
  replace condition 2\textsuperscript{o} by the hypothesis that there exists
  an $\mathcal{O}_X$-module ample for $f$
  (\hyperref[II.4.6.11]{4.6.11}).
  \item[\rm{(ii)}] Let $Y$ be a quasi-compact scheme such that there exists
  an ample $\mathcal{O}_Y$-module. For a $Y$-scheme $X$ to be projective, it
  is necessary and sufficient that it be $Y$-isomorphic to a closed
  sub-$Y$-scheme of a projective bundle of the form $\mathbf{P}_Y^r$. The
  condition is evidently sufficient.

\oldpage[II]{105}

  Conversely, if $X$ is projective over $Y$, it is quasi-projective, hence
  there exists a $Y$-immersion $j$ of $X$ into some $\mathbf{P}_Y^r$
  (\hyperref[II.5.3.3]{5.3.3}), which is closed by
  (\hyperref[II.5.4.4]{5.4.4}) and
  (\hyperref[II.5.5.3]{5.5.3}).
  \item[\rm{(iii)}] The argument of
  (\hyperref[II.5.5.3]{5.5.3}) shows that, for every prescheme $Y$ and
  every integer $r\geq0$, the structural morphism $\mathbf{P}_Y^r\to Y$ is
  surjective, for if we set
  $\mathcal{S}=\mathbf{S}_{\mathcal{O}_Y}(\mathcal{O}_Y^{r+1})$, we
  evidently have $\mathcal{S}_y=\mathbf{S}_{k(y)}((k(y))^{r+1})$
  (\hyperref[II.1.7.3]{1.7.3}), hence $(\mathcal{S}_n)_y\neq0$ for every
  $y\in Y$ and every $n\geq0$.
  \item[\rm{(iv)}] It follows from the examples of Nagata \cite{II-26} that
  there are proper morphisms which are not quasi-projective.
\end{enumerate}

\begin{proposition}[5.5.5]
\label{II.5.5.5}
\begin{enumerate}
  \item[\rm{(i)}] A closed immersion is a projective morphism.
  \item[\rm{(ii)}] If $f:X\to Y$ and $g:Y\to Z$ are projective
  morphisms, and if $Z$ is a quasi-compact scheme or a prescheme whose
  underlying space is Noetherian, then $g\circ f$ is projective.
  \item[\rm{(iii)}] If $f:X\to Y$ is a projective $S$-morphism, then
  $f_{(S')}:X_{(S')}\to Y_{(S')}$ is projective for every extension
  $S'\to S$ of the base prescheme.
  \item[\rm{(iv)}] If $f:X\to Y$ and $g:X'\to Y'$ are projective
  $S$-morphisms, then so is $f\times_Sg$.
  \item[\rm{(v)}] If $g\circ f$ is a projective morphism and if $g$ is
  separated, then $f$ is projective.
  \item[\rm{(vi)}] If $f$ is projective, then so is
  $f_{\mathrm{red}}$.
\end{enumerate}
\end{proposition}

(i) follows immediately from (\hyperref[II.3.1.7]{3.1.7}). Here we must
prove (iii) and (iv) separately, because of the restriction imposed on
$Z$ in (ii) (cf. (\hyperref[I.3.5.1]{I, 3.5.1})). To prove (iii), we
reduce to the case where $S=Y$ (\hyperref[I.3.3.11]{I, 3.3.11}), and the
assertion then follows immediately from (\hyperref[II.5.5.1]{5.5.1, b)})
and from (\hyperref[II.3.5.3]{3.5.3}). To prove (iv), we are immediately
reduced to the case where $X=\mathbf{P}(\mathcal{E})$,
$X'=\mathbf{P}(\mathcal{E}')$, $\mathcal{E}$ (resp. $\mathcal{E}'$) being a
quasi-coherent $\mathcal{O}_Y$-module (resp. $\mathcal{O}_{Y'}$-module) of
finite type. Let $p$, $p'$ be the canonical projections of $T=Y\times_SY'$
to $Y$ and $Y'$ respectively; by (\hyperref[II.4.1.3.1]{4.1.3.1}), we have
$\mathbf{P}(p^*(\mathcal{E}))=\mathbf{P}(\mathcal{E})\times_YT$,
$\mathbf{P}(p'^*(\mathcal{E}'))=\mathbf{P}(\mathcal{E}')\times_{Y'}T$, whence
\begin{align*}
\mathbf{P}(p^*(\mathcal{E}))\times_T\mathbf{P}(p'^*(\mathcal{E}'))
&=(\mathbf{P}(\mathcal{E})\times_YT)\times_T
(T\times_{Y'}\mathbf{P}(\mathcal{E}'))\\
&=\mathbf{P}(\mathcal{E})\times_Y
(T\times_{Y'}\mathbf{P}(\mathcal{E}'))
=\mathbf{P}(\mathcal{E})\times_S\mathbf{P}(\mathcal{E}')
\end{align*}
on replacing $T$ by $Y\times_SY'$ and using
(\hyperref[I.3.3.9.1]{I, 3.3.9.1}). Now, $p^*(\mathcal{E})$ and
$p'^*(\mathcal{E}')$ are of finite type over $T$
(\hyperref[0.5.2.4]{0, 5.2.4}), hence so is
$p^*(\mathcal{E})\otimes_{\mathcal{O}_T}p'^*(\mathcal{E}')$; since
$\mathbf{P}(p^*(\mathcal{E}))\times_T\mathbf{P}(p'^*(\mathcal{E}'))$
is identified with a closed sub-prescheme of
$\mathbf{P}(p^*(\mathcal{E})\otimes_{\mathcal{O}_T}p'^*(\mathcal{E}'))$
(\hyperref[II.4.3.3]{4.3.3}), this completes the proof of (iv). To obtain
(v) and (vi), we can apply (\hyperref[I.5.5.13]{I, 5.5.13}), since every
closed sub-prescheme of a projective $Y$-scheme is a projective $Y$-scheme by
(\hyperref[II.5.5.1]{5.5.1, a)}).

It remains to prove (ii); by virtue of the hypothesis on $Z$, this follows
from (\hyperref[II.5.5.3]{5.5.3}), from
(\hyperref[II.5.3.4]{5.3.4, (ii)}) and
(\hyperref[II.5.4.2]{5.4.2, (ii)}).

\begin{proposition}[5.5.6]
\label{II.5.5.6}
If $X$ and $X'$ are two projective $Y$-schemes, then $X\amalg X'$ is a
projective $Y$-scheme.
\end{proposition}

This is an evident consequence of (\hyperref[II.5.5.2]{5.5.2}) and
(\hyperref[II.4.3.6]{4.3.6}).

\begin{proposition}[5.5.7]
\label{II.5.5.7}
Let $X$ be a projective $Y$-scheme and $\mathcal{L}$ a $Y$-ample
$\mathcal{O}_X$-module; for every section $f$ of $\mathcal{L}$ over $X$,
$X_f$ is affine over $Y$.
\end{proposition}

\oldpage[II]{106}

The question being local on $Y$, we may suppose
$Y=\operatorname{Spec}(A)$; moreover
$X_{f^{\otimes n}}=X_f$, so by replacing $\mathcal{L}$ by a suitable
$\mathcal{L}^{\otimes n}$, we may suppose $\mathcal{L}$ very
ample for the structural morphism $q:X\to Y$
(\hyperref[II.4.6.11]{4.6.11}). The canonical homomorphism
$\sigma:q^*(q_*(\mathcal{L}))\to\mathcal{L}$ is therefore surjective and
the corresponding morphism
\[
r=r_{\mathcal{L},\sigma}:X\to P=\mathbf{P}(q_*(\mathcal{L}))
\]
is an immersion such that $\mathcal{L}=r^*(\mathcal{O}_P(1))$
(\hyperref[II.4.4.4]{4.4.4}); moreover, since $X$ is proper over $Y$,
the immersion $r$ is closed (\hyperref[II.5.4.4]{5.4.4}). Furthermore, we
have by definition $f\in\Gamma(Y,q_*(\mathcal{L}))$, and $\sigma^\flat$ is
the identity of $q_*(\mathcal{L})$; it then follows from formula
(\hyperref[II.3.7.3.1]{3.7.3.1}) that $X_f=r^{-1}(D_+(f))$; $X_f$
is consequently a closed sub-prescheme of the affine scheme $D_+(f)$ and is
therefore indeed an affine scheme.

When in particular we take $Y=X$, we obtain (taking into account
(\hyperref[II.4.6.13]{4.6.13, (i)})) the following corollary, whose
direct proof is moreover immediate:

\begin{corollary}[5.5.8]
\label{II.5.5.8}
Let $X$ be a prescheme and $\mathcal{L}$ an invertible $\mathcal{O}_X$-module.
For every section $f$ of $\mathcal{L}$ over $X$, $X_f$ is affine
over $X$ (and consequently an affine scheme when $X$ is an affine scheme).
\end{corollary}

\subsection{Chow's lemma.}
\label{subsection:II.5.6}

\begin{theorem}[5.6.1]
\label{II.5.6.1}
(Chow's lemma). Let $S$ be a prescheme and $X$ an $S$-scheme of finite type.
Suppose that one of the following hypotheses is satisfied:
\begin{enumerate}
  \item[\rm{a)}] $S$ is Noetherian.
  \item[\rm{b)}] $S$ is a quasi-compact scheme, and $X$ has finitely many
  irreducible components.
\end{enumerate}
Under these conditions:
\begin{enumerate}
  \item[\rm{(i)}] There exist an $S$-scheme $X'$ quasi-projective over $S$
  and an $S$-morphism $f:X'\to X$ which is projective and surjective.
  \item[\rm{(ii)}] We can take $X'$ and $f$ such that there exists an open
  subset $U\subset X$ for which $U'=f^{-1}(U)$ is dense in $X'$ and the
  restriction of $f$ to $U'$ is an isomorphism $U'\simeq U$.
  \item[\rm{(iii)}] If $X$ is reduced (resp. irreducible, integral), we can
  suppose $X'$ reduced (resp. irreducible, integral).
\end{enumerate}
\end{theorem}

The proof proceeds in several steps.

A) We can first reduce to the case where $X$ is irreducible. Indeed, under
hypothesis a), $X$ is Noetherian, so that, under either hypothesis, the
irreducible components $X_i$ of $X$ are finite in number. If the theorem is
proved for each of the reduced closed sub-preschemes of $X$ having the $X_i$
as underlying spaces, and if $X'_i$ and $f_i:X'_i\to X_i$ are the prescheme
and the morphism corresponding to $X_i$, then the prescheme $X'$, the sum of
the $X'_i$, and the morphism $f:X'\to X$ whose restriction to each $X'_i$ is
$j_i\circ f_i$ ($j_i$ the canonical injection $X_i\to X$) meet the
requirements. It is immediate indeed that $X'$ is reduced if the $X'_i$ are;
on the other hand, we satisfy (ii) by taking for $U$ the union of the sets
$U_i\cap(\bigcup_{j\neq i}X_j)^{\complement}$. Finally, since the $X'_i$ are
quasi-projective over $S$, so is $X'$

\oldpage[II]{107}

(\hyperref[II.5.3.6]{5.3.6}); likewise, the morphisms $X'_i\to X$ are
projective by (\hyperref[II.5.5.5]{5.5.5, (i) and (ii)}), hence $f$ is
projective (\hyperref[II.5.5.6]{5.5.6}), and it is evidently surjective by
definition.

B) Now suppose $X$ irreducible. Since the structural morphism
$r:X\to S$ is of finite type, there is a finite covering $(S_i)$ of $S$ by
affine open subsets, and for each $i$ a finite covering $(T_{ij})$ of
$r^{-1}(S_i)$ by affine open subsets, the morphism $T_{ij}\to S_i$ being
affine and of finite type, hence quasi-projective
(\hyperref[II.5.3.4]{5.3.4, (i)}); since under each of the hypotheses a),
b) the immersion $S_i\to S$ is quasi-compact, it is a quasi-projective
morphism (\hyperref[II.5.3.4]{5.3.4, (ii)}), hence the restriction of
$r$ to $T_{ij}$ is a quasi-projective morphism
(\hyperref[II.5.3.4]{5.3.4, (iii)}). Denote the $T_{ij}$ by $U_k$
$(1\leq k\leq n)$. For each index $k$ there exists an open immersion
$\varphi_k:U_k\to P_k$, where $P_k$ is projective over $S$
(\hyperref[II.5.3.2]{5.3.2} and \hyperref[II.5.5.2]{5.5.2}). Let
$U=\bigcap_kU_k$; since $X$ is irreducible and the $U_k$ are non-empty, $U$ is
non-empty, and consequently everywhere dense in $X$; the restrictions to $U$
of the $\varphi_k$ define a morphism
\[
\varphi:U\to P=P_1\times_SP_2\times_S\cdots\times_SP_n
\]
such that the diagrams
\[
\tag{5.6.1.1}
\label{II.5.6.1.1}
\xymatrix{
U\ar[r]^\varphi\ar[d]_{j_k} & P\ar[d]^{p_k}\\
U_k\ar[r]_{\varphi_k} & P_k
}
\]
are commutative, $j_k$ being the canonical injection $U\to U_k$ and $p_k$ the
canonical projection $P\to P_k$. If $j$ is the canonical injection $U\to X$,
the morphism $\psi=(j,\varphi)_S:U\to X\times_SP$ is an immersion
(\hyperref[I.5.3.14]{I, 5.3.14}). Under hypothesis a), $X\times_SP$ is
locally Noetherian (\hyperref[II.3.4.1]{3.4.1} and
\hyperref[I.6.3.7]{I, 6.3.7 and 6.3.8}); under hypothesis b),
$X\times_SP$ is a quasi-compact scheme
(\hyperref[I.5.5.1]{I, 5.5.1 and 6.6.4}); in both cases the closure
$X'$ in $X\times_SP$ of the sub-prescheme $Z$ associated to $\psi$ (and whose
underlying space is $\psi(U)$) exists, and $\psi$ factors as
\[
\tag{5.6.1.2}
\label{II.5.6.1.2}
\psi:U\xrightarrow{\psi'}X'\xrightarrow{h}X\times_SP
\]
where $\psi'$ is an open immersion and $h$ a closed immersion
(\hyperref[I.9.5.10]{I, 9.5.10}). Let
$q_1:X\times_SP\to X$ and $q_2:X\times_SP\to P$ be the canonical projections;
we set
\[
\tag{5.6.1.3}
\label{II.5.6.1.3}
f:X'\xrightarrow{h}X\times_SP\xrightarrow{q_1}X
\]
and
\[
\tag{5.6.1.4}
\label{II.5.6.1.4}
g:X'\xrightarrow{h}X\times_SP\xrightarrow{q_2}P.
\]
We will see that $X'$ and $f$ meet the requirements.

C) Let us first show that $f$ is projective and surjective, and that the
restriction of $f$ to $U'=f^{-1}(U)$ is an isomorphism of $U'$ onto $U$.
Since the $P_k$ are projective over $S$, so is $P$
(\hyperref[II.5.5.5]{5.5.5, (iv)}), hence $X\times_SP$ is projective over
$X$ (\hyperref[II.5.5.5]{5.5.5, (iii)}), and so is $X'$, which is a
closed sub-prescheme of $X\times_SP$. On the other hand, we have
$f\circ\psi'=q_1\circ(h\circ\psi')=q_1\circ\psi=j$, hence $f(X')$ contains
the everywhere dense open set $U$ of $X$; but $f$ is a closed morphism
(\hyperref[II.5.5.3]{5.5.3}), hence $f(X')=X$. Let us now note that
$q_1^{-1}(U)=U\times_SP$ is induced on an open subset of $X\times_SP$, and
by definition the prescheme $U'=h^{-1}(U\times_SP)$ is induced by $X'$ on
the open subset $U'$; it is consequently the closure with

\oldpage[II]{108}

respect to $U\times_SP$ of the prescheme $Z$
(\hyperref[I.9.5.8]{I, 9.5.8}). Now the immersion $\psi$ factors as
\[
U\xrightarrow{\Gamma_\varphi}U\times_SP\xrightarrow{j\times1}X\times_SP,
\]
and since $P$ is separated over $S$, the graph morphism $\Gamma_\varphi$ is
a closed immersion (\hyperref[I.5.4.3]{I, 5.4.3}), hence $Z$ is a
closed sub-prescheme of $U\times_SP$, from which we conclude that $U'=Z$.
Since moreover $\psi$ is an immersion, the restriction of $f$ to $U'$ is an
isomorphism onto $U$, inverse to $\psi'$; finally, by definition of $X'$,
$U'$ is dense in $X'$.

D) Let us now prove that $g$ is an immersion, which will establish that $X'$
is quasi-projective over $S$, since $P$ is projective over $S$. Set
\begin{align*}
V_k&=\varphi_k(U_k) &&\text{(open subset of $P_k$)},\\
W_k&=p_k^{-1}(V_k) &&\text{(open subset of $P$)},\\
U'_k&=f^{-1}(U_k),\qquad U''_k=g^{-1}(W_k)
&&\text{(open subsets of $X'$)}.
\end{align*}
It is clear that the $U'_k$ form an open covering of $X'$; we will first see
that the same holds for the $U''_k$, by proving that
$U'_k\subset U''_k$. For this it will suffice to show that the diagram
\[
\tag{5.6.1.5}
\label{II.5.6.1.5}
\xymatrix{
U'_k\ar[r]^{g|U'_k}\ar[d]_{f|U'_k} & P\ar[d]^{p_k}\\
U_k\ar[r]_{\varphi_k} & P_k
}
\]
is commutative. Now the prescheme $U'_k=h^{-1}(U_k\times_SP)$ is induced by
$X'$ on the open subset $U'_k$, and is therefore the closure of
$Z=U'\subset U'_k$ with respect to $U'_k$ (\hyperref[I.9.5.8]{I, 9.5.8}).
To show the commutativity of diagram (\hyperref[II.5.6.1.5]{5.6.1.5}), it
therefore suffices, since $P_k$ is an $S$-scheme, to show that on composing
this diagram with the canonical injection $U'\to U'_k$ (or, what amounts to
the same by virtue of the isomorphism of $U'$ and $U$, with $\psi$), we
obtain a commutative diagram (\hyperref[I.9.5.6]{I, 9.5.6}). But by
definition the diagram thus obtained is none other than
(\hyperref[II.5.6.1.1]{5.6.1.1}), whence our assertion.

The $W_k$ therefore form an open covering of $g(X')$; to prove that
$g$ is an immersion, it suffices to see that each of the restrictions
$g|U''_k$ is an immersion into $W_k$
(\hyperref[I.4.2.4]{I, 4.2.4}). For this, consider the morphism
$u_k:W_k\xrightarrow{p_k}V_k\xrightarrow{\varphi_k^{-1}}U_k\to X$; since
$X$ is separated over $S$, the graph morphism
$\Gamma_{u_k}:W_k\to X\times_SW_k$ is a closed immersion
(\hyperref[I.5.4.3]{I, 5.4.3}), hence the graph
$T_k=\Gamma_{u_k}(W_k)$ is a closed sub-prescheme of $X\times_SW_k$; if
we show that this sub-prescheme dominates $U'$, it will also dominate the
sub-prescheme induced by $X'$ on the open subset $U''_k$ of $X'$, by
(\hyperref[I.9.5.8]{I, 9.5.8}). Since the restriction of $q_2$ to $T_k$
is an isomorphism onto $W_k$, the restriction of $g$ to $U''_k$ will be an
immersion into $W_k$, and our assertion will be proved. Let $v_k$ be the
canonical injection $U'\to X\times_SW_k$; we must show that there exists a
morphism $w_k:U'\to W_k$ such that $v_k=\Gamma_{u_k}\circ w_k$. By
definition of the product, it suffices to prove that
$q_1\circ v_k=u_k\circ q_2\circ v_k$
(\hyperref[I.3.2.1]{I, 3.2.1}), or, on composing on the right

\oldpage[II]{109}

with the isomorphism $\psi':U\to U'$, that
$q_1\circ\psi=u_k\circ q_2\circ\psi$. But since $q_1\circ\psi=j$ and
$q_2\circ\psi=\varphi$, our assertion follows from the commutativity of
diagram (\hyperref[II.5.6.1.1]{5.6.1.1}), taking into account the
definition of $u_k$.

E) It is clear that, since $U$, and consequently $U'$, is irreducible, so is
$X'$ in the preceding construction, and the morphism $f$ is therefore
birational (\hyperref[I.2.2.9]{I, 2.2.9}). If moreover $X$ is reduced, so
is $U'$, hence $X'$ is reduced
(\hyperref[I.9.5.9]{I, 9.5.9}). This completes the proof.

\begin{corollary}[5.6.2]
\label{II.5.6.2}
Suppose that one of the hypotheses a), b) of
(\hyperref[II.5.6.1]{5.6.1}) is satisfied. For $X$ to be proper over $S$,
it is necessary and sufficient that there exist a scheme $X'$ projective over
$S$ and a surjective $S$-morphism $f:X'\to X$ (which is then projective by
(\hyperref[II.5.5.5]{5.5.5, (v)})). When this is the case, we can
moreover choose $f$ so that there exists a dense open subset $U$ of $X$ such
that the restriction of $f$ to $f^{-1}(U)$ is an isomorphism
$f^{-1}(U)\simeq U$ and $f^{-1}(U)$ is dense in $X'$. If moreover $X$
is irreducible (resp. reduced), we can suppose that so is $X'$; when $X$
and $X'$ are irreducible, $f$ is a birational morphism.
\end{corollary}

The condition is sufficient by virtue of (\hyperref[II.5.5.3]{5.5.3}) and
(\hyperref[II.5.4.3]{5.4.3, (ii)}). It is necessary, for with the
notation of (\hyperref[II.5.6.1]{5.6.1}), if $X$ is proper over $S$, then
$X'$ is proper over $S$, since it is projective over $X$, hence proper over
$X$ (\hyperref[II.5.5.3]{5.5.3}), and our assertion follows from
(\hyperref[II.5.4.2]{5.4.2, (ii)}); moreover, since $X'$ is
quasi-projective over $S$, it is projective over $S$ by virtue of
(\hyperref[II.5.5.3]{5.5.3}).

\begin{corollary}[5.6.3]
\label{II.5.6.3}
Let $S$ be a locally Noetherian prescheme, $X$ an $S$-scheme of finite type
over $S$, and $f_0:X\to S$ its structural morphism. For $X$ to be proper
over $S$, it is necessary and sufficient that for every morphism of finite
type $S'\to S$, $(f_0)_{(S')}:X_{(S')}\to S'$ be a closed morphism. It even
suffices that this condition be satisfied for every $S$-prescheme of the form
$S'=S\otimes_{\mathbf Z}\mathbf Z[T_1,\ldots,T_n]$ ($T_i$ indeterminates).
\end{corollary}

The condition being evidently necessary, let us show that it is sufficient.
The question being local on $S$ and $S'$ (\hyperref[II.5.4.1]{5.4.1}), we
may suppose $S$ and $S'$ affine and Noetherian. By virtue of Chow's lemma,
there exist a projective $S$-scheme $P$, an immersion $j:X'\to P$ and a
projective and surjective morphism $f:X'\to X$ such that the diagram
\[
\xymatrix{
X\ar[d]_{f_0} & X'\ar[l]_f\ar[d]^j\\
S & P\ar[l]_r
}
\]
is commutative. Since $P$ is of finite type over $S$, the first hypothesis
implies that the projection $q_2:X\times_SP\to P$ is a closed morphism. But
the immersion $j$ factors as $j=q_2\circ(f\times1)\circ\Gamma_j$,
where $\Gamma_j:X'\to X'\times_SP$ is the graph morphism, which is a
closed immersion since $P$ is separated over $S$
(\hyperref[I.5.4.3]{I, 5.4.3}). Now $f$, being projective, is proper
(\hyperref[II.5.5.3]{5.5.3}), hence $f\times1$ is closed. We conclude that
$j$ is a closed immersion, and is consequently proper
(\hyperref[II.5.4.2]{5.4.2, (i)}). Moreover, the structural morphism
$r:P\to S$ is projective, hence proper (\hyperref[II.5.5.3]{5.5.3}), hence
$f_0\circ f=r\circ j$ is proper
(\hyperref[II.5.4.2]{5.4.2, (ii)}); finally, since $f$ is surjective,
$f_0$ is proper by virtue of (\hyperref[II.5.4.3]{5.4.3}).

To prove the proposition using only the second hypothesis, it suffices to
show that it implies the first. Now, if $S'$ is affine and of finite type
over $S=\operatorname{Spec}(A)$,

\oldpage[II]{110}

we have $S'=\operatorname{Spec}(A[c_1,\ldots,c_n])$
(\hyperref[I.6.3.3]{I, 6.3.3}), and $S'$ is therefore isomorphic to a
closed sub-prescheme of $S''=\operatorname{Spec}(A[T_1,\ldots,T_n])$ ($T_i$
indeterminates). In the commutative diagram
\[
\xymatrix{
X\times_SS'\ar[r]^{1_X\times j}\ar[d]_{(f_0)_{(S')}} &
X\times_SS''\ar[d]^{(f_0)_{(S'')}}\\
S'\ar[r]_j & S''
}
\]
$j$ and $1_X\times j$ are closed immersions
(\hyperref[I.4.3.1]{I, 4.3.1}) and $(f_0)_{(S'')}$ is closed by
hypothesis; hence so is $(f_0)_{(S')}$.

